\section{无穷级数 III}

\begin{theorem}
    设正项无穷数列 $\{a_n\}$ 满足
    $$
    \frac{a_n}{a_{n + 1}} = 1 + \frac{1}{n} + \frac{\beta}{n \ln n} + o\ab(\frac{1}{n \ln n})
    $$
    那么：
    \begin{itemize}
        \item 当 $\beta > 1$ 时，$\sum\limits_{n = 1}^{\infty} a_n$ 收敛；
        \item 当 $\beta < 1$ 时 $\sum\limits_{n = 1}^{\infty} a_n$ 发散。
    \end{itemize}

    \begin{proof}
        将 $a_n$ 与 $\frac{1}{n (\ln n)^{p}}$ 比较，使用比值判别法。
    \end{proof}
\end{theorem}

\subsection{级数和积分的关系}

\begin{theorem}
    对于一个非负函数 $f(x)$，若积分 $\displaystyle I = \int_a^{+ \infty} f(x) \dd{x}$ 收敛，那么对于任意序列 $\{A_n\}$ 满足：$a = a_1 < a_2 < \dots$ 且 $\lim\limits_{n \to \infty} a_n = +\infty$，都有：
    $$
    \sum_{n = 1}^{\infty} \int_{A_n}^{A_{n + 1}} f(x) \dd{x} = I
    $$

    \begin{proof}
        使用广义积分的定义和夹逼定理即可。
    \end{proof}
\end{theorem}

\subsection{一般正项级数的敛散性判别法}

\begin{definition}[交错级数]
    设数列 $\{x_n\}$ 满足 $x_n = u_n (-1)^n\ (u_n > 0)$，那么级数 $\sum_{i = 1}^{\infty} x_n$ 称为\textbf{交错级数}。若满足 $u_n$ 递减且 $\lim\limits_{n \to \infty} u_n = 0$，那么称其为\textbf{Lebniz 级数}。 
\end{definition}



\begin{lemma}[分部求和]
    设 $\{a_n\}, \{b_n\}$ 是两个数列，那么对于任意 $n \in \mathbb{N}^+$ 有：
    $$
    \sum_{k = 1}^{n} a_k b_k = \sum_{k = 1}^{n - 1} S_k (b_k - b_{k + 1}) + S_n b_n
    $$
    其中 $S_k = \sum\limits_{i = 1}^{k} a_i$。
\end{lemma}

\begin{lemma}[Abel 引理]
    设 $\{a_n\}, \{b_n\}$ 是两个数列，$S_k = \sum\limits_{i = 1}^{k} a_i$，若满足：
    \begin{itemize}
        \item $\{b_n\}$ 单调；
        \item $|A_k|$ 存在上界 $M$。
    \end{itemize}
    那么
    $$
    \forall\,p \in \mathbb{N}^+: \abs{\sum_{k = 1}^{p} a_k b_k} \le M\ab(\abs{b_1} + 2 \abs{b_p})
    $$
    \begin{proof}
        不妨设 $\{b_n\}$ 单调递减，那么：
        $$
        \begin{aligned}
            \abs{\sum_{k = 1}^{p} a_k b_k} & = \abs{S_p b_p + \sum_{k = 1}^{p - 1} S_k (b_k - b_{k + 1})} \\
            & \le \abs{S_p} \abs{b_p} + \sum_{k = 1}^{p - 1} \abs{A_k} (b_k - b_{k + 1}) \\
            & \le M \abs{b_p} + M \sum_{k = 1}^{p - 1} (b_k - b_{k + 1}) \\
            & = M\ab(\abs{b_p} + b_1 - b_p) \le M \ab(\abs{b_1} + 2 \abs{b_p})
        \end{aligned}
        $$
    \end{proof}
\end{lemma}

由此我们可以得到 Dirichlet 判别法。

\begin{theorem}[Dirichlet 判别法]
    设 $\{a_n\}, \{b_n\}$ 是两个数列，$S_n = \sum\limits_{i = 1}^{n} a_i$，若：
    \begin{itemize}
        \item $\{b_n\}$ 单调且 $\lim\limits_{n \to \infty} b_n = 0$；
        \item $\abs{S_n}$ 有上界 $M$。
    \end{itemize}
    那么级数 $\sum\limits_{n = 1}^{\infty} a_n b_n$ 收敛。

    \begin{proof}
        考虑 $k = n + 1, n + 2, \dots, n + p$，那么根据 Abel 引理：
        $$
        \abs{\sum_{k = n + 1}^{n + p} a_k b_k} \le 2M \ab(\abs{y_{n + 1}} - 2 \abs{y_{n + p}})
        $$
        而显然 $n \to \infty$ 是 $\rhs \to 0$，于是根据 Cauchy 收敛法则得证。
    \end{proof}
\end{theorem}

此外更改 Dirichlet 判别法的条件，可以得到 Abel 判别法：

\begin{theorem}[Abel 判别法]
    设 $\{a_n\}, \{b_n\}$ 是两个数列，$S_n = \sum\limits_{i = 1}^{n} a_i$，若：
    \begin{itemize}
        \item $\{b_n\}$ 单调有界；
        \item $\sum_{i = 1}^{\infty} a_n$ 收敛。
    \end{itemize}
    那么级数 $\sum\limits_{n = 1}^{\infty} a_n b_n$ 收敛。

\end{theorem}